
\section{TODO Explainability criteria}

Explainability in programs is not linked to a specific characteristic and cannot be easily achieved. Programmers often don't understand their own code after some time, although it (hopefully) even has documentation. If programs, which are to be understood by the developer, are now to be created automatically, much more attention must be paid to clarity so that the viewer is not overwhelmed with information.

In the following, we discuss what the criteria for explainability can be and how they can be taken into account in this approach. For this purpose, a measure of explainability must be found above all. In order to optimize in this respect, the existing possibilities of the GP for generating explainable programs are explored and in addition, some possible and implemented extensions of GP for even more explainability are proposed and discussed.


\subsubsection{TODO Possibilities to maximize explainability with native settings}

Explanability can already be improved by an adapted parameterization. This includes, among others:

- strong typing. Automated type-casts seem very unintuitive
- small function set
- small complexity as second optimisation

\subsection{TODO Measuring explainability}


The essential goal of this approach is therefore to find a measure of the explainability in programs that can be optimized in addition to fitness. This measure must apply equally to all program components and types of problems.

The selected complexity measures usually depend on the node [cite], the depth of the tree [cite] and variations thereof [todo, evaluation of func, ...].

More than anything else, the actual sense is mainly based on naturally shaped contexts. 

The developer probably asks himself most of all what he has to change in his program to improve it even more. Thereby as little as possible should be changed. Our complexity measure has therefore been based primarily on the following question: ``How much must be changed in the reference program to achieve a certain fitness?

TODO alle Neuerungen zeigen

\begin{itemize}
    \item Tree size: tree length, tree depth, tree width
    \item pattern and structures.
    \item function types.
    \item variable types:
    \item encapsulation
    \item possibility to recognize relevant parts. 
    \item level, on which variable
    \item encapsulation of the variable.
    \item order of logic: 
    \item encapsulation of isolated complexity:
    \item if then else:
    \item encapsulation of program parts.
    \item reuse of identical structures within a solution.
    \item obviousness of meaning. 
\end{itemize}

\paragraph{Explainability distance requirements}
The measure of explainability has to meet a number of requirements. \label{xaiMeasurementRequirements}

TODO Darstellung als (1) Prämisse 1, ...
\begin{itemize}
    \item Depending on the perceived explainability
    \item Fulfilling the prerequisites for a distance measure
        \begin{itemize}
        \item $d(x,y) \ge 0$  non-negativity
        \item $d(x,y) = 0 \Leftrightarrow x = y$  identity of indiscernibles
        \item $d(x,y) = d(y,x)$  symmetry
        \item $d(x,y) \le d(x,z) + d(z,y)$  triangle inequality
        \end{itemize}
    \item Fine graduation of dimensions.
    \item Univariate or low-dimansional multivariate formula.
    \item No fundamental limits the possibilities of the GP occur.
\end{itemize}

The fact that the measure is a distance is a fundamental requirement. If a measure were chosen which is not based on these premises, no ranking of the programs could be generated. Therefore the distance measure must be able to show even small differences. In order to be able to comprehend the explanatory measure, not too many factors should influence the value. A multivariate formula that includes a multitude of factors is conceivable here, but since the weighting of possible influencing factors is very complex and that of the artificially constructed unit ``explicability'' is very arbitrary and each factor must be justified and tested in detail, we try to keep the formula as simple as possible.  
Furthermore, the genetic programming should not be significantly restricted in its possibilities. Explainable programs are a nice thing, but they are of no use if the behavior of the genetic programming is restricted to such an extent that good candidates can no longer be found at all or only with great difficulty.

\paragraph{Absolute and relative factors}
The factors of the complexity measure can be selected either absolute or relative to the reference tree.  Absolute properties are, for example, whether a particular function is part of a program or how many nodes a tree has. Relative factors refer to changes compared to the reference program.  Both aspects have advantages and disadvantages. Absolute values are equally valid for all trees and thus provide a more general measure. Relative factors have the advantage, however, that they support the use of already known structures from the reference tree more strongly. The developer is also better shown where there is still potential for improvement in his program. However, a program that achieves the performance of the reference program in a much easier way would be harder to find. However, since we can assume that the expert has chosen his core program sensibly, relative factors should be an essential part of the explanatory measure.

\paragraph{Optimizing GP candidates with fitness and complexity}
It would also be possible to calculate the fitness of programs not only on the basis of their performance, but to choose a combination of complexity and performance. For example, it would be conceivable to combine the percentage increase in fitness with the increase in complexity. However, there is also the problem that the fitness of the programs would be misappropriated if programs with completely different results had the same fitness. Instead of introducing another artificially created measure, the extent of which would still have to be clarified for the GP, the candidates are arranged on a Pareto front. This front contains all Pareto efficient solutions, so each entry is either fitter or more explainable than another. The developer is then free to choose at his own discretion the programs that offer the best combination for him.

\subsection{Relevant factors for explainability}
% todo, do i need a lot of sources here?
%todo graphics for everything

In this part, most of the factors influencing explainability are explained and it is discussed whether they are suitable as factors for a distance measure. 

\subsubsection{todo In use GP complexity measurements}

The longer an expression is, the more complex it is and the easier it is to lose track of what is actually happening. Candidates should therefore be kept as small as possible and the size of a tree should be part of the final measurement. 

\paragraph{TODO node count \label{eq:expl:nodeCount}}
The number of nodes is probably the most obvious factor for the complexity of trees. The more nodes a tree has, the more confusing its display and the more complex its contents become. The number of nodes records even small changes and can be used both as an absolute value and as a relative value if the changes are counted towards the reference tree. In addition, this unit of measurement can also be easily extended, for example, by weighting nodes.

\paragraph{Measurable values: Tree depth and width  \label{eq:expl:treeDepth}}
Another widespread form of complexity is the depth of a tree. The depth is often a good indicator, because the number of nodes increases exponentially downward, and the complexity caused by the stringing together of functions also increases. Furthermore, in genetic programming, trees are often created based on a depth that must not be exceeded. However, if trees are not created in full mode, so that any branch can stop growing at any time, the depth of a tree is not only a weak attempt to approximate the number of nodes as a measure, but also a misleading indicator.
Moreover, the depth of a tree does not satisfy the triangle inequality, since trees that have the same program can have different depths.
As discussed in ref [todo refmodularity], deep trees are actually preferable to wide trees for the sake of explanation.

If a tree has a very large average of nodes in the same plane, the high number of branches can create more complex relationships than without. [todo show,( -, +, *, 1, 2, 3, 4)] The resulting complexity can be reduced by the increased use of the grow method. Thus, although the trees are easier to understand, complex functions are harder to achieve. In order not to limit the possibilities of genetic programming too much, the 'full' method was not artificially limited.

\subsubsection{Encapsulation, Divide and conquer}
Problems can be understood most easily if they are broken down into the smallest possible parts that can be analyzed separately. An extreme counterexample is a neural network, where it is hardly possible to look at parts in isolation to see their function for the overall result. The more such subdivisions there are, the easier it is to derive a logical function from a part of the program. Creating encapsulated programs does not only deliver better explainability, improvements in subtrees can also be achieved more easily, because the components can be exchanged modularly.

\paragraph{Encapsulation: Distribution and Dimensionality of the Variable}

The variable has a special meaning in encapsulation, because the program tries to place only the operators around the variable in such a way that the correct result is obtained.
If input variables do not occur in a fit program, its influence on the problem is small or their information content can be taken also from other variables. The dimensionality of the problem could thus be limited without great risk. The more isolated the variables are from each other, the less complex their interaction is.

For example, if only one variable or a small subset of variables is in use within a sub-tree, the real dependence on variables may not be reduced, but the interaction can be recognized more precisely and it is not necessary to switch between the different functions of the variables.

The occurrence of variables could be weighted separately with a relatively high value to force a reduced use of them. This was not part of the distance measure in the tests performed due to keeping it simple.

\paragraph{Enclosure. Complex program parts} 
Under certain circumstances a program cannot avoid very complex partial calculations. Here it can make a difference whether this part is calculated at a single point or whether the connection results from many different combinations.

\paragraph{TODO Encapsulation: if then else}
The conditional if function plays a special role here. Although it is the most complex function with 3 input parameters, it creates more structure than all other operators because its child trees are not combined but can be considered separately. Nevertheless, the subtrees must have a relation to each other in order to harmonize together. This combination creates encapsulated programs in a clear way that humans can understand more easily. It is always to be welcomed when these are contained in programs, but this should not be artificially forced, as they are not useful at every point.

For the complexity measure, the If function could be given a low weighting, but the possibility of evaluating only the most complex subtree should also be considered. For the sake of simplicity, however, the use of If loops in the program is not yet part of the complexity measure.

\subsubsection{Recognition: Known Patterns and Structures}

The human brain is very good at recognizing patterns. This ability is deeply rooted in us through evolution, as strange faces or predators can be quickly recognized as potential threats \cite{humanPatternRecognition}. In fact, this ability has become so pronounced in evolution that faces are often seen in places where there are none, like spotting the image of Jesus on a toast.

\paragraph{TODO Function types}
The patterns of some functions are easier to understand than others because they represent more natural relationships. For example, the ``+'' function requires two input values, double as much as the $tan$ function, but its result is easier to understand. If some mathematical relationships should not be considered because they are not expected, special mathematical functions such as $tan, sin, cos, e^x$ should be excluded from genetic programming.

Normally, functions that find unfamiliar correlations will rarely occur due to the low rate of successful crossover candidates. In this way, bad functions are already used less frequently through selection.

In addition, operators can be weighted to indicate their occurrence by a lower explanatory power. Giving the functions weights depends on the developer and most improvements could be achieved when complex functions are not used at all.

\paragraph{TODO Constants}
Constants are actually not an important factor for explainability. Nevertheless, there are also differences between them. For example, constants are easier to process the fewer numbers they have. Therefore, they should have as few decimal places as necessary. Nevertheless, efficient fine-tuning of constants should not be prevented. Constants were therefore rounded to an accuracy of $0.005$.
Hote that when solving primarily filigree, mathematical problems, rounding should be avoided completely.

Besides, if a constant is found to be similar to a mathematical constant, for example $3.14$, a pattern is clearly recognized by any mathematician. Thus, it can be replaced by its real constant for both aiming for better performance and also showing the developer patterns that are often dedicated to a specific range of problems. Examples of such constants are: $e, \sqrt{2}, \frac{1}{2}, \frac{1}{3}, \frac{1}{4}, \frac{2}{3}, ...)$ 

\paragraph{TODO Pattern: Order}
All roads lead to Rome. Tests, sequences, and mathematical formulas do not necessarily have to follow a certain order. Even the position of a node does not allow you to conclude with certainty that it is relevant. Nevertheless, certain procedures can be recognized more easily. This is the case, for example, if a calculation that runs differently due to commutativity or associativity could still mix the most complex parts first.
[TODO examples. $2*a+5$ is not $+5+a*2$. Showing more known structures?]

As already mentioned, the variables should be affected as late as possible in a tree, to mix and complexify their context too early.

\paragraph{TODO Pattern: Reappearing structures within a solution}
If certain calculations are repeated in a tree, it is easier to recognize a pattern. By crossover between trees, the same branches are automatically exchanged over and over again, which naturally increases the probability that structures appear repeatedly. Compared to normal GP, an increase in the proportion of crossover in the evolution would allow the same patterns to be inserted again to an even greater extent. However, this was neither necessary nor should it affect the performance of the GP process.

A much more suitable way to have the same function appear more often is to avoid inverse functions or similar functions when selecting the operators to generate new trees. These include, for example, the comparison operators $<, <=, >, >=, ==, !=$, angle functions $sin, cos, tan$, or Boolean operators such as $And, Or, Xand, Xor, Xand, Nand, Nor, Not$. All these operators can be generated from subsets without much trouble, so some of them can be omitted.


\subsubsection{TODO (paragraph?) Recognition of origin, number and type of changes}

In our case, the expert's reference program is regarded as a special case of recognition. The logical core structure it provides is not only a gene already recognized by the developer (at least hopefully) which can therefore be easily adopted in any candidate - if its structure is recognized by the developer, it also provides an easy to understand anchor point for the rest of the code. The reference to the origin should therefore be relatively easy to establish.

A very lightweight solution for this would be to add the tree as a candidate so often in each generation that clear descendants are created again and again through crossover. If, however, too many of the same candidate trees are added to a population, the diversity and innovation generated by new candidates is inhibited.

Measuring the difference is the main component of our measurement. To achieve this, number of changes required to modify the reference tree to the found solution is computed, which is a known problem in theoretical computer science, called the tree edit distance [todo ref]

Other options that were discussed before deciding to use the tree edit distance were:
\begin{itemize}
\item number of positions to be applied

From an explanatory point of view, it is actually interesting to count the number of branches inserted. As a form of encapsulation, a tree with a small number of additional branches would be preferable.
However, it is out of question as the main measure. On the one hand, an expansion at several places may be unavoidable, which would be the next problem: It does not allow for simple gradations of complexity.

In order to offer interfaces to the genetic programming where there is a lot of potential for improvement, the expert should also formulate and structure the initial program in such a way that extensions at as many points as possible are understandable and can be integrated meaningfully. In the case of Mountain Car, the ``simple'' program was extended in such a way that Action 0 can be integrated without much effort.

\item Arbitrary generated function.

For each node that is not part of the original, its arity is added as penalty.
This measure has been tested among others, but compared to the actual number of changes, it is only an artificially created aid that does not offer as results with the same precision and plausibility as the tree edit distance.

\item Accumulated number of changes since start

This measurement is unsuitable for two reasons. First, it makes the use of crossover impossible, as it is unclear which parents complexity value should be accumulated. Second, as a distance measure it does not satisfy the triangle inequality.
$test$ [todo proves triangle inequality]

\end{itemize}

\paragraph{TODO weighting functions}

TODO reference?
Weighting functions is very practical for an expert as he can add his own preferences for functions to the measurement. This is enabled by not only computing the distance itself but also analysing the required mapping of performed edits to change the origin tree to the found solution. 

For For example, the developer could consider the sum-operation to be less complicated than the exponent-function. This would mean, that the user considers sums to be more understandable. This extra measurement is very personalised though and thus was not evaluated specifically for this thesis.
